% META: {"id": "1SPE-VARALEA-EV-B-corrige", "chapitre": "1SPE-VARIABLES-ALEATOIRES", "type_objet": "corrige_evaluation", "evaluation_ref": "1SPE-VARALEA-EV-B", "version": "B", "capacites_codes": ["C1","C2","C3","C4","C5","C6","C7"], "sources_inspiration": [], "status": "generated"}
% BEGIN-VERIFY
% from sympy import Rational, sqrt
% # Ex 1
% E = 2*Rational(1,2)+5*Rational(3,10)+8*Rational(1,5)
% assert E == Rational(41,10)
% E2 = 4*Rational(1,2)+25*Rational(3,10)+64*Rational(1,5)
% assert E2 == Rational(223,10)
% V = E2 - E**2; assert V == Rational(549,100)
% E_Y = 2*E - 3; assert E_Y == Rational(82-30,10); assert E_Y == Rational(52,10)
% # Ex 2
% probabilities = [Rational(81,256), Rational(108,256), Rational(54,256), Rational(12,256), Rational(1,256)]
% P0 = probabilities[0]; P1 = probabilities[1]
% assert P0 == Rational(81,256) and P1 == Rational(27,64)
% E_X = sum(k*probabilities[k] for k in range(5)); assert E_X == 1
% V_X = sum(k*k*probabilities[k] for k in range(5))-E_X**2
% assert V_X == Rational(3,4)
% Pge1 = 1-P0; assert Pge1 == Rational(175,256)
% # Ex 3
% E_B = 80 - 2000*Rational(1,50); assert E_B == 40
% B_500 = 500*40; assert B_500 == 20000
% # Ex 4 : simulation et fluctuation
% mu_4 = Rational(25,10); sigma_4 = Rational(18,10)
% seuil_4 = 2*sigma_4/sqrt(81)
% assert seuil_4 == Rational(4,10)
% assert Rational(29,10) - mu_4 == Rational(4,10)
% assert (Rational(29,10) - mu_4) == seuil_4
% assert Rational(381,400) == Rational(9525,10000)
% assert Rational(381,400)*100 == Rational(1905,20)
% END-VERIFY

\begin{corrige}{1SPE-VARALEA-EV-B}

%===============================================================
\section*{Exercice 1 \hfill (8 points) — C1, C2, C4}
%===============================================================

\textbf{Question 1 — Tableau de loi.}

\[
\begin{array}{|c|c|c|c|}
\hline
x_i & 2 & 5 & 8 \\
\hline
P(X=x_i) & \frac{5}{10} = \frac{1}{2} & \frac{3}{10} & \frac{2}{10} = \frac{1}{5} \\
\hline
\end{array}
\]

\medskip
\textbf{Question 2 — Espérance.}

$E(X) = 2 \times \frac{1}{2} + 5 \times \frac{3}{10} + 8 \times \frac{1}{5} = 1 + \frac{3}{2} + \frac{8}{5} = \frac{10+15+16}{10} = \frac{41}{10} = 4{,}1$.

En moyenne, on obtient $4{,}1$ points par tirage.

\medskip
\textbf{Question 3 — Variance et écart type.}

$E(X^2) = 4 \times \frac{1}{2} + 25 \times \frac{3}{10} + 64 \times \frac{1}{5} = 2 + \frac{15}{2} + \frac{64}{5} = \frac{20+75+128}{10} = \frac{223}{10}$.

$V(X) = \frac{223}{10} - \left(\frac{41}{10}\right)^2 = \frac{2230-1681}{100} = \frac{549}{100} = 5{,}49$.

$\sigma(X) = \sqrt{5{,}49} \approx 2{,}34$.

\medskip
\textbf{Question 4.}

$E(Y) = 2E(X) - 3 = 2 \times 4{,}1 - 3 = 5{,}2$.

%===============================================================
\section*{Exercice 2 \hfill (7 points) — C2, C3}
%===============================================================

\textbf{Question 1.} L'arbre comporte quatre niveaux. À chaque nœud, la branche « réussi » porte $\frac14$ et la branche « manqué » $\frac34$.

\medskip
\textbf{Question 2.} $P(X=0)=\left(\frac34\right)^4=\frac{81}{256}$.

\medskip
\textbf{Question 3.} Chacun des quatre chemins comportant exactement une réussite a pour probabilité $\frac14\left(\frac34\right)^3=\frac{27}{256}$. Donc $P(X=1)=4\times\frac{27}{256}=\frac{27}{64}$.

\medskip
\textbf{Question 4.}
\[
\begin{array}{|c|c|c|c|c|c|}\hline
x&0&1&2&3&4\\ \hline
P(X=x)&\frac{81}{256}&\frac{108}{256}&\frac{54}{256}&\frac{12}{256}&\frac1{256}\\ \hline
\end{array}
\]
On obtient $E(X)=1$, $E(X^2)=\frac74$, donc $V(X)=\frac34$ et $\sigma(X)=\frac{\sqrt3}{2}$.

\medskip
\textbf{Question 5.} $P(X\geq1)=1-P(X=0)=1-\frac{81}{256}=\frac{175}{256}$.

%===============================================================
\section*{Exercice 3 \hfill (5 points) — C5}
%===============================================================

\textbf{Question 1.} Pas de sinistre : $B = 80$. Sinistre : $B = 80 - 2000 = -1920$.

\medskip
\textbf{Question 2.}

\[
\begin{array}{|c|c|c|}
\hline
b_i & -1920 & 80 \\
\hline
P(B=b_i) & \frac{1}{50} & \frac{49}{50} \\
\hline
\end{array}
\]

\medskip
\textbf{Question 3.} $E(B) = 80 - 2000 \times \frac{1}{50} = 80 - 40 = 40$~euros.

$E(B) > 0$ : le contrat est rentable pour l'assureur.

\medskip
\textbf{Question 4.} Bénéfice total moyen : $500 \times 40 = 20\,000$~euros.

%===============================================================
\section*{Exercice 4 \hfill (5 points) — C6, C7}
%===============================================================

\baremeIndicatif{Q1: 1 pt; Q2: 1 pt; Q3: 1,5 pt; Q4: 1,5 pt}

\bigskip
\textbf{Question 1 — Ligne Python.}

\commentaireMarge{La division porte sur la taille de l'échantillon.}

\begin{python}
    return somme / n
\end{python}

\medskip
\textbf{Question 2.} $\sqrt{81} = 9$, donc
\[
\frac{2\sigma}{\sqrt{n}} = \frac{2 \times 1{,}8}{9} = \frac{3{,}6}{9} = 0{,}4.
\]

\medskip
\textbf{Question 3.} $|m - \mu| = |2{,}9 - 2{,}5| = 0{,}4$.

\commentaireMarge{Cas limite : l'inégalité est large, l'égalité convient.}

L'écart vaut exactement $0{,}4$, c'est-à-dire le seuil lui-même. L'inégalité
demandée étant large, la condition $|m - \mu| \leq \dfrac{2\sigma}{\sqrt{n}}$
est \textbf{vérifiée}. Cet échantillon est donc compté parmi ceux de
l'intervalle.

\medskip
\textbf{Question 4.} La proportion vaut
\[
\frac{381}{400} = 0{,}9525 = 95{,}25\,\%.
\]
Elle est cohérente avec l'ordre de grandeur attendu, environ $95\,\%$.

\end{corrige}
